% Complete replacement of the released reciprocal-source subsection.
% Reviewed source: Version 2, SHA256 38d458ba4178083766c5ae56a894da595fe27170c1ef07f6b49750fc5edf7c37.
% Selected source: H19 / high40 / beta12 / anchor25 / scalar32 / 29 full + 8 beta.
% No endpoint count, source pin, certificate hash, or coefficient claim is supplied here.
\subsection{The selective reciprocal comparison source}
\label{subsec:quartic-source}
We use the curve, cyclic group of integer indices, pole predicate and
arithmetic inputs of Section~\ref{sec:degree8-original}. The present
source uses a sharper consequence of the same height construction,
thirty-seven direct coordinate contacts, and two kinds of reciprocal
comparison. These arithmetic changes are proved here before the
quadratic lift. They are not justified by a change of source digest.
Throughout,
\[
 K=\Q(\sqrt2),\quad d=3-\sqrt2,\quad D=d^2=11-6\sqrt2,
 \qquad E:y^2=x^3-Dx,\qquad G=2E(K)=\Z P_*.
\]
An element of $K$ is represented by its two rational coefficients.
Multiplication and inversion have their literal meanings
\[
 (a_0,a_1)(b_0,b_1)=(a_0b_0+2a_1b_1,a_0b_1+a_1b_0),\qquad
 a^{-1}=\frac{(a_0,-a_1)}{a_0^2-2a_1^2}\quad(a\ne0).
\]
The denominator is nonzero by the irrationality of $\sqrt2$.
All affine subtraction charts and their nonzero-abscissa guards remain.

% BLOCK HEIGHT19.
\begin{theorem}[A sharper five-contact height estimate]\label{thm:quartic-height19}
Keep $S_*$, $\mathcal D$ and $M_5$ from Theorem~\ref{hyp:o8-height}, and put
\[
 B_*=\prod_{p\le76176}p,\qquad
 A_{19}=2^{109}210^{1800},\qquad H_{19}=2A_{19}B_*^{19}+128.
\]
For every $a\in\mathbb Q\setminus\mathcal D$,
\[
 h(a)\le H_{19}M_5(a)^{19},
 \qquad \log_2H_{19}<2093966.
\]
Here $h$ is the absolute logarithmic Weil height.
\end{theorem}
\begin{proof}
We use the geometric construction and compatible representations in the
proof of Theorem~\ref{hyp:o8-height}, not merely its final exponent-$24$
inequality. Write $R=B_*M_5(a)$ and $C=210^{300}$.
The geometric comparison there is
\[
 h(a)\le2\max(0,h_F(A))+128.
\]
All $h_F$ in this proof use the original stable Faltings metric.
The CM estimate in that proof is $h_F(A)<854R$.

In the non-CM case the same construction gives an absolutely simple
quaternionic surface $A/L$, a normal field $L'/\mathbb Q$, and the final
primitive weight-two newform $f$, after the finite twist has been
absorbed, with
\[
 [L:\mathbb Q]\le576,\quad d'=[L':\mathbb Q]\le2304R,\quad
 N\le CR^2,\quad \operatorname{rad}(N)\mid R.
\]
Its two-dimensional compatible representation $r_\ell$ satisfies
\[
 V_\ell(A)|_{G_{L'}}\simeq
       (r_\ell\oplus r_\ell)|_{G_{L'}},\qquad
 (\det r_\ell)^2|_{G_{L'}}=\chi_{\rm cyc}^2|_{G_{L'}}.
\]
These are (LR3)--(LR5) and the determinant calculation in the existing
modularity argument. There is no bounded-degree hypothesis on the
coefficient field of $f$.

Let $\epsilon$ be the nebentypus of this final form. In the Tate-module
convention, with arithmetic Frobenius, $\det r_\ell=\epsilon\chi_{\rm cyc}$;
see \cite[\S12.5, (12.5.1)--(12.5.2) and Remark 12.5.1,
pp.~119--120]{diamond-im}. Equivalently, at a good prime its
characteristic polynomial is $X^2-a_pX+p\epsilon(p)$.
A consistent contragredient convention replaces the finite character by
its inverse, with the same kernel and order. The character $\epsilon^2$
is trivial on $G_{L'}$, so its image order divides $d'$. Consequently
\[
 t=\operatorname{ord}(\epsilon)\le2d'\le4608R.
\]
This bounds the character order, not the coefficient-field degree.

Put $H=\ker\epsilon\subset(\mathbb Z/N\mathbb Z)^*$ and
\[
 \Gamma_H=\left\{
 \begin{pmatrix}a&b\\c&d\end{pmatrix}\in\Gamma_0(N):
 d\bmod N\in H\right\}.
\]
Since $\epsilon(-1)=1$, this group contains $-I$.
Let $X_H=X_1(N)/H$ for the diamond action and $J_H=\operatorname{Jac}(X_H)$.
The curve, diamond action and quotient are defined over $\mathbb Q$.
The rational model here parametrizes a point of exact order $N$, up to
$H$, and is the model specified by \cite[\S2.1, (2.1.8)]{mascot-moduli};
it introduces no cyclotomic field-degree factor.

Write $q:J_1(N)\to A_f$ for the newform quotient. Its pulled-back
differentials are spanned by all Galois conjugates of $f$; see
\cite[Theorem~1(i),(iv), p.~526]{shimura-factors} and
\cite[\S3.1]{baker-et-al-finiteness}. Every $h\in H$ fixes those differentials,
since its eigenvalue on each conjugate is a conjugate of $\epsilon(h)=1$;
the diamond convention is \cite[\S3.2.1]{baker-et-al-finiteness}.
Homomorphisms of abelian varieties in characteristic zero are determined
by their differentials. Hence $q\circ\langle h\rangle=q$.
For
\[
 e_H=\frac1{|H|}\sum_{h\in H}\langle h\rangle
       \in\operatorname{End}_{\mathbb Q}^0(J_1(N))
\]
we therefore have $qe_H=q$.

Let $H_{\rm eff}$ be the image of $H$ in
$\operatorname{Aut}(X_1(N))$, and let $\pi:X_1(N)\to X_H$ be the quotient.
On Jacobians,
\[
 \pi^*\pi_*=\sum_{h\in H_{\rm eff}}h,\qquad
 \pi_*\pi^*=[|H_{\rm eff}|].
\]
The averages over $H$ and $H_{\rm eff}$ agree. In particular the trivial
action of $-1$ is not counted twice in the first identity.
The image of $e_H$ is thus $\mathbb Q$-isogenous to $J_H$, and its
restriction of $q$ is surjective. Poincare reducibility makes $A_f$
a $\mathbb Q$-isogeny factor of $J_H$.

The accepted realization of $r_\ell$ in $A_f$ puts it in $V_\ell(J_H)$.
Over $L'$, Faltings' Hom theorem and extension of scalars give a nonzero
homomorphism $A\to J_H$. Absolute simplicity and Poincare reducibility
give an $L'$-isogeny $A\times Z\to J_H$. In particular
$g=\dim J_H\ge2$. The Jacobian has its canonical principal polarization;
no bound on a polarization degree of $A$ or $Z$ is being assumed.

The map $X_H\to X(1)$ has degree
\[
 d_H=[\operatorname{SL}_2(\mathbb Z):\Gamma_H]
     =Nt\prod_{p\mid N}(1+1/p).
\]
There is no factor of two here, because $-I\in\Gamma_H$.
It is branched only over the cusp and the two elliptic points, so
$\deg_B(X_H)\le d_H$. The congruence-curve genus bound gives
$g\le1+d_H/12\le d_H$, using $g\ge2$.
These degree and genus conventions, and the resulting height inequality,
are recorded in \cite[\S4, (4.1)--(4.2), p.~8]{vonKaenel2021}.
For every prime $p\ge5$, $(1+1/p)^6\le p$, while the possible factors
at two and three have product at most two. It follows that
\[
 d_H\le9216CR^{19/6}<T,\qquad
 T=2^{14}CR^{19/6},\qquad g<T.
\]
By \cite[Theorem~1.1.1 and \S2.3]{javanpeykar},
\[
 h_F(J_H)\le13000000\,d_H^5g<2^{24}T^6.
\]
This is the generic Belyi-curve estimate; no squarefree-level restriction
or trivial-nebentypus assumption enters it.

Apply \cite[Theorem~1.9(1), pp.~7--8]{GaudronRemond2023} with source
$(J_H)_{L'}$ and target $A\times Z$. Combining its degree bound with
half-logarithmic isogeny-height variation gives
\[
 h_F(A)+h_F(Z)\le h_F(J_H)+V,\quad
 V=2g^2\left(8g^2\log(7g)+\log d'
     +\log\max(1,\log d',h_F(J_H))\right).
\]
The Bost lower bound in \cite[(2.2), p.~5]{vonKaenel2021} gives
$h_F(Z)>-3g/2$. Put $u=\log T$. Then
\[
 V\le16T^4(2+u)+2T^2(25+7u)\le T^6.
\]
Indeed $\log d'<8+\log R\le8+u$, and the last logarithm in $V$ is
less than $17+6u$. For $T\ge16$ the inequalities
$2+\log T\le T^2/32$ and $25+7\log T\le T^4/4$ follow at $16$
and remain true by differentiation. Thus
\[
 h_F(A)<(2^{24}+1)T^6+\tfrac32T
       <2^{25}T^6=A_{19}R^{19}.
\]
The earlier CM estimate is also less than this bound.
The geometric comparison and $M_5(a)\ge1$ prove the assertion.
Finally $B_*<2^{109445}$, $210<2^8$ and $A_{19}B_*^{19}>64$ give
\[
 H_{19}<4A_{19}B_*^{19}
       <2^{\,2+109+1800\cdot8+19\cdot109445}
       =2^{2093966}.
\]
\end{proof}

% BLOCK ROOT37: no denominator-clearing variable is added to this source.
\begin{lemma}[Thirty-seven contacts for rational coefficients]
\label{lem:quartic-rational37}
Keep $K=\Q(\sqrt2)$, $d=3-\sqrt2$, $D=d^2$, and
$G=2E(K)=\Z P_*$ from Section~\ref{sec:degree8-original}.
Let $A=[n]P_*,S=[s]P_*,T=[t]P_*,Z=[z]P_*$ be nonidentity points.
Order the eight rational coefficients of
$x(A),x(S),x(T),x(Z)$ as $a_1,\ldots,a_8$. Put
\[
 R_{\rm c}=47027,\quad W_0=25878772920,\quad c=18484837801,
 \quad V_0=32W_0,
 \qquad A_i=R_{\rm c}a_i+3,\quad w_i=A_iv-u.
\]
Use the following thirty-seven nonzero raw contacts:
\begin{equation}\label{eq:quartic-contacts37}
 r_0=v/V_0,\qquad u-\delta v\quad(\delta\in\mathcal D),
 \qquad -(w_i-\delta v)\quad(1\le i\le8,\ \delta\in\mathcal D).
\end{equation}
Here and below the order within $\mathcal D$ is
$(0,1,2401,-343/32)$. The common five positions are numbered $0,\ldots,4$;
the other positions are numbered $5,\ldots,36$ in the displayed order.

Every normalized odd contact of $u/v$ or $w_i/v$ outside $S_*$ is
detected by an odd valuation at a raw position belonging to that slope.
If the radicals for both slopes $u/v,w_i/v$ have their $e$th powers
dividing an integer $N\ne0$, then, with $H=H_{19}$,
\[
 h(a_i)<2H|N|^{19/e}+14.
\]
Consequently the two coefficients of any one of the four abscissas give
\begin{equation}\label{eq:quartic-direct-height}
 h(x(Q))<4H|N|^{19/e}+30<7H|N|^{19/e}
 \quad(Q\in\{A,S,T,Z\})
\end{equation}
whenever all its required positions have those divisibilities.

Conversely, for any fixed eight rational coefficients, the inherited
ordinary-prime construction supplies these raw contacts with every finite
contact a positive selected prime and $r_0$ a positive rational square
times a selected prime. Repeated positions may have equal values.
The same pole-completeness construction applies to these values.
\end{lemma}
\begin{proof}
Nonzero raw contacts imply $v\ne0$ and exclude the four finite branch
values for every slope. At a prime outside $S_*$ all differences of
distinct members of $\mathcal D$ are units. Write $m$ for the minimum
of the valuations of the two coordinates of one slope. If an odd finite
normalized contact has $m$ even, its distinguished raw linear form has
odd valuation. If $m$ is odd, another finite form has valuation $m$.
For an odd contact at infinity use the denominator in the first case
and a finite form in the second. The unit $V_0$ and the signs do not
change parity. This proof uses rational coordinates and also covers
negative valuations; it does not presuppose that a contact is prime.

Apply Theorem~\ref{thm:quartic-height19} to the two slopes. The identity
$R_{\rm c}a_i+3=u/v+w_i/v$ and the elementary logarithmic height inequalities
give the first bound, since $\log(6R_{\rm c})<14$. For
$x(Q)=a+b\sqrt2$ use
$h(x(Q))\le h(a)+h(b)+h(\sqrt2)+\log2$.
This proves \eqref{eq:quartic-direct-height}, since $H>128$ and $|N|\ge1$.

For the converse choose a positive integer $D_0$ with $D_0a_i\in\Z$
for every $i$, after fixing the eight coefficients. This integer is a
choice in the existence proof, not a variable of the source. Set
\[
 T_0=V_0D_0^2,\qquad p_0=c+W_0y,\qquad
 u=c+W_0x,\qquad v=T_0p_0.
\]
For each distinct slope
$b\in\mathcal D\cup\bigcup_i(A_i-\mathcal D)$ put
$K_b=T_0b\in W_0\Z$ and $\Psi_b=u-K_bp_0$.
Together with $p_0$, these are at most thirty-seven integer affine forms,
each congruent to $c$ modulo $W_0$. Distinct slopes give independent
linear parts. Exact duplicate slopes are removed only while applying
the prime theorem and restored at all their original positions afterward.

The value and successor admissibility in
Section~\ref{subsec:o8-pattern-proof} remains valid. At primes dividing
$W_0$ use the same residue $c$. If a prime $\ell>29$ divides $D_0$,
all $K_b$ vanish modulo $\ell$, and $u=p_0=1$ avoids zero and minus one.
If $\ell>29$ does not divide $D_0$, the change from $(u,p_0)$ to $(u,v)$
is invertible. For $\ell=31,37,41$ there are at most eight slope residues,
since $R_{\rm c}a_i\equiv0$; for $\ell\ge43$ there are at most thirty-six.
Choose distinct missing residues $b_0,e_0$ and order them so that
$e_0-b_0\ne-T_0^{-1}$. Then
\[
 v=(e_0-b_0)^{-1},\qquad u=b_0v
\]
make every $\Psi_b$ different from zero and minus one, and make
$p_0=v/T_0$ different from zero and minus one. Interchanging the two
missing residues permits the required ordering, because $\ell$ is odd.

Choose a fixed positive cone
$u/v>\max_i|A_i|+2413$, $v>0$. Enlarge the fixed initial cutoff past
$82$, every prime dividing $D_0$, and the nonzero determinants of the
distinct forms. The inherited prime-pattern proof then applies with its
same unused third coordinate and with complexity at most $35$.
For its fiber upper estimates, fixing $\Psi_{b_0}=p$ gives the same
one-variable forms
$\Psi_b=p+(K_{b_0}-K_b)p_0$.
All coefficients, the cone, the denominator $D_0$ and the cutoff are
fixed before the box grows. Thus the proof uses precisely the inherited
fixed-system prime and successor estimates, not a uniform estimate in a
growing denominator. It gives the positive ordinary primes and the
same-prime elliptic-order conditions used there.

Now $r_0=D_0^2p_0$. At a prime $\ell\mid D_0$ outside $S_*$, both $u$
and $p_0$ are units, $v_\ell(v)=2v_\ell(D_0)$, and
$v_\ell(A_i)\ge-v_\ell(D_0)$. Hence $w_i=A_iv-u$ is a unit.
Both slopes have negative valuation $-2v_\ell(D_0)$, and their normalized
infinite depths are even. No new normalized odd contact is created at
such a denominator prime. At $p_0$, the odd valuation of $r_0$ is one.

For each occurrence choose its private comparison witnesses sufficiently
deep that its norm root is a unit at its assigned prime. The remaining
denominator valuations in the pole argument occur in the square of that
root. The raw contacts are positive, their reduced denominators are
squares, and they are squares in $\Q_2,\Q_3,\Q_{17}$ by the retained
residue $c$ and the square factor $D_0^2$. Lemma~\ref{hyp:o8-pole}
therefore supplies the pole witnesses. Equal prime occurrences use the
same orientation and may use separate private points. No condition on
reduction orders at different primes has been added.
\end{proof}

% BLOCK RECIPROCAL: the released complete proofs with the stated uniform range.
\begin{lemma}[Reciprocal logarithm and its error]
\label{lem:quartic-reciprocal-log}
\label{lem:quartic-uniform-reciprocal}
Put $f=1/x$, $z=-x/y$, and define
\[
 a_0=1,\qquad a_1=-\frac25,\qquad
 a_j=\frac{4j-5}{4j+1}a_{j-1}\quad(j\ge2),\qquad
 G(w)=\sum_{j\ge0}a_jw^j.
\]
On the formal branch at the identity,
\[
 \log_E(z)=zG(Df^2).
\]
For every $J\ge0$ put $G_J(w)=\sum_{j=0}^{J}a_jw^j$.
For every integer $10\le e\le64$, set $m=\lceil e/4\rceil$ and
\[
 G_{m-1}(w)=\sum_{j=0}^{m-1}a_jw^j,\qquad
 L_e(z,f)=zG_{m-1}(Df^2).
\]
Let $p>76176$ be split and choose either place above it. If the point is
nonidentity and lies in the formal group at that place, so
$r=v_p(z)\ge1$ and $v_p(f)>0$, then
\[
 v_p\bigl(\log_E(z)-L_e(z,f)\bigr)\ge(4m+1)r,
 \qquad \frac{L_e(z,f)}{\log_E(z)}\in1+p^e\Z_p.
\]
\end{lemma}
\begin{proof}
The curve equation gives
\[
 f=z^2(1-Df^2),\qquad w=Df^2,\qquad Dz^4=\frac{w}{(1-w)^2}.
\]
Coefficient comparison in the defining recurrence gives
\[
 4w(1-w)G'(w)+(1+w)G(w)=1-w.
\]
Indeed the coefficient of $w$ is $5a_1+a_0=-1$, and each later
coefficient gives $(4j+1)a_j=(4j-5)a_{j-1}$. Differentiating the
curve identity yields $z\,dw/dz=4w(1-w)/(1+w)$. Consequently
\[
 \frac{d}{dz}\{zG(w)\}=\frac{1-w}{1+w}
                   =(1+4Dz^4)^{-1/2}.
\]
The square root has constant term one. This is the normalized invariant
differential for $z$, and $zG(w)$ has zero constant term, which proves
the formal identity.

At the stated place, $D$ is a unit and the formal branch gives
$v_p(f)=2r$, hence $v_p(w)=4r$. The denominator of $a_j$, for $j\ge1$,
divides $5\prod_{k=2}^j(4k+1)$ and therefore divides $(4j+1)!$.
Thus
\[
 v_p(a_j)\ge-v_p((4j+1)!)\ge-\frac{4j+1}{p-1}.
\]
For $j\ge m$ with $4j+1<p$, the coefficient is $p$-integral and
$v_p(a_jw^j)\ge4mr$. Otherwise $j\ge19044$, so $3j\ge4m$ and
$(4j+1)/(p-1)\le j$. Therefore
\[
 v_p(a_jw^j)\ge4jr-j\ge4mr.
\]
These lower bounds tend to infinity. They control the entire infinite
tail, including terms whose denominators are divisible by $p$.
Multiplication by $z$ proves the first estimate. Formal logarithm isometry
gives $v_p(\log_E(z))=r$, and proves the relative estimate.
\end{proof}

\begin{lemma}[Global nonvanishing]
\label{lem:quartic-reciprocal-nonzero}
For every $f\in K$ and every $J\ge2$,
\[
                 G_J(Df^2)\ne0.
\]
In particular $L_e(z,f)\ne0$ whenever $z\ne0$ and $10\le e\le64$.
\end{lemma}
\begin{proof}
Let $\nu_2$ be the rational exponent-of-two valuation and put
\[
 v(a+b\sqrt2)=\tfrac12\nu_2(a^2-2b^2)
             =\min\{\nu_2(a),\nu_2(b)+\tfrac12\}.
\]
The two terms in the norm have valuations of opposite parity when finite.
Norm multiplicativity and the displayed minimum prove the valuation
identities. Its finite values lie in $\tfrac12\Z$, and $v(D)=0$.
Every $a_j$ with $j\ge1$ has valuation one, because $a_1=-2/5$ and
all subsequent recurrence factors have odd numerator and denominator.
If $f=0$, the assertion is immediate. If $t=v(f)\ge0$, the constant
term of $G_J(Df^2)$ is uniquely least in valuation. If $t<0$, then
$t\le-1/2$ and the nonconstant valuations $1+2jt$ strictly decrease;
the last is at most $1-J<0$, below the constant term as well. In either
case a uniquely least term proves nonvanishing. No good-reduction
assumption at two is used.
\end{proof}

% BLOCK SELECTIVE_SPECIFICATION: the selected beta-only source.
\subsubsection{Selective comparison systems}
For this source put
\begin{equation}\label{eq:quartic-selected-parameters}
 H=H_{19},\qquad \kappa=40,\qquad q=12,
 \qquad E_{\rm sc}=32,\qquad A_{\rm an}=25.
\end{equation}
Here $\kappa$ is a precision, while $K$ continues to denote the field.
Retain $X^2=\beta-2$, the proxy equation \eqref{eq:o8-Lucas}, all four
nonidentity root points $A,S,T,Z$, the two common auxiliary points,
and the base pole and norm predicates. In the scalar batch use orders
$(1,1,1,1,32,32)$ on the two coefficients each of
$f(J_0),f(S),z(S)/z(J_0)$. In the anchor use
$c_0=3$, $c_{i+1}=c_i^2$ for $0\le i<25$ and $C=c_{25}$.
Thus $C$ is computed by witnesses; it is not a new coefficient of the
polynomial. Replace the denominator-cleared height targets by precisely
the eight coefficients in Lemma~\ref{lem:quartic-rational37}.

Here is a uniform specification of the two comparator types. Write
$n_1=n,n_2=s,n_3=t,n_4=z$ and
$(b_1,b_2,b_3,b_4)=(\beta,\xi,\tau,\zeta)$; put
$Q_0=P_*$ and $Q_i=[n_i]P_*$. For a chosen subset
$I\subset\{1,2,3,4\}$ introduce private nonidentity primary points
$P_0=[h]P_*$ and $P_i=[m_i]P_*$ for $i\in I$. Use the base affine
subtraction charts for
\[
 R_0=P_0-[4]Q_0,\qquad R_i=P_i-[4]Q_i\quad(i\in I).
\]
Every old subtraction guard is retained: in particular the two input
abscissas are different and these differences are nonidentity. The
shared $J=[j]P_*\ne O$ and its reciprocal parameters are also retained.
No identity-allowing difference model is used in this specification.

For precisions $e_0$ and $e_i\le e_0$ define
\[
 C_i=L_{e_i}(P_i)/L_{e_0}(P_0).
\]
The ordered field tests are first
\[
 f(P_0),\ (f(P_i))_{i\in I},\ f(J)^\sigma,
 \ \phi_0,\ (\phi_i)_{i\in I},
 \qquad \phi_i=f(R_i)^\sigma,
\]
all of order one, followed by $(C_i-b_i)_{i\in I}$ of orders $e_i$,
and then $z(R_0)^\sigma/z(J)^\sigma$ of order $e_0$ and
$(z(R_i)^\sigma/z(J)^\sigma)_{i\in I}$ of orders $e_i$.
Here $\sigma(\sqrt2)=-\sqrt2$, and every list indexed by $I$ is in
increasing order. The already-conjugated common inputs are used once,
not conjugated a second time.

Before these paired field tests, retain the single rational test
$\varepsilon=\theta^2-2$ of order one. A field expression
$a+b\sqrt2$ of order $e$ gives the two rational tests
$\varepsilon b,a+b\theta$, both of order $e$.
Each resulting list has its own constant-left norm tree and the same
pole predicate as the base compiler. The two choices are
\[
\begin{array}{c|c|c|c}
 \text{type}&I&e_0&(e_i)_{i\in I}\\ \hline
 \text{full}&\{1,2,3,4\}&40&(12,40,40,40)\\
 \text{beta only}&\{1\}&12&(12).
\end{array}
\]
Their exact rational order lists are, respectively,
\begin{align*}
 &\bigl(1^{[23]},12^{[2]},40^{[8]},12^{[2]},40^{[6]}\bigr),\\
 &\bigl(1^{[11]},12^{[6]}\bigr).
\end{align*}
where $e^{[r]}$ means $r$ consecutive copies of $e$.
These have lengths $41,17$; the lists specify expressions as well as
orders. In particular the beta-only denominator has precision $12$.

Assign these types to the thirty-seven positions as follows:
\begin{equation}\label{eq:quartic-mode-partition}
\begin{array}{c|c|c}
 \text{positions}&\text{role}&\text{comparison}\\ \hline
 0\text{--}4&\text{common}&\text{full}\\
 5\text{--}12&x(A)&\text{beta only}\\
 13\text{--}20&x(S)&\text{full}\\
 21\text{--}28&x(T)&\text{full}\\
 29\text{--}36&x(Z)&\text{full}.
\end{array}
\end{equation}
There are $29$ full and $8$ beta-only positions. These are
source call counts, not polynomial variable counts. No raw position is
removed or merged when two slopes or primes coincide.

\begin{lemma}[The local conclusions of the two comparators]
\label{lem:quartic-selective-local}
At every $p>76176$ where a raw contact has odd valuation, a full
comparator gives
\[
 \beta\in\Z_p,\quad \beta\equiv n\pmod {p^{12}},\qquad
 (\xi,\tau,\zeta)\in\Z_p^3,\quad
 (\xi,\tau,\zeta)\equiv(s,t,z)\pmod {p^{40}}.
\]
A beta-only comparator gives the displayed beta conclusion, with no
Lucas conclusion asserted at that position.
No ordinary-prime hypothesis is needed in this direction.
\end{lemma}
\begin{proof}
Let $R=v_p(r)$ be the odd valuation of the raw contact. The all-left
constant path of the norm tree and Lemma~\ref{hyp:o8-pole} exclude
$R<0$. For $R>0$, descent through that tree gives, for a rational test
$Y$ assigned order $e$,
\[
 v_p(Y)>(e-1)R\ge e-1,\qquad v_p(Y)\ge e.
\]
This argument is unchanged for each of the two tree sizes and for
every positive integer order. It does not assume a prime raw value.

The root-error test makes $\theta$ a $p$-adic unit and selects a root
$\rho$ of $2$ with $\theta-\rho\in p\Z_p$ and $\theta+\rho$ a unit.
Thus $p$ splits. Since
$\theta-\rho=\varepsilon/(\theta+\rho)$, the paired inequalities
$v_p(\varepsilon b),v_p(a+b\theta)\ge e$ imply
$v_p(a+b\rho)\ge e$, even when $b$ is not integral.
The order-one $f$ tests put all primary points in the formal group at
this place, and $J,R_0,(R_i)$ in the formal group at its conjugate.

Formal-logarithm homomorphism and isometry
\cite[IV.6.4 and VII.2.1--2.2]{silverman} give
\[
 v_p((h-4)/j)\ge e_0,\qquad
 v_p((m_i-4n_i)/j)\ge e_i.
\]
Here $j$ and all indices are integers and $j\ne0$. Therefore
$h\equiv4\pmod {p^{e_0}}$, $m_i\equiv4n_i\pmod {p^{e_i}}$,
and $h$ is a unit. The same homomorphism at the primary place gives
$\log_E(P_i)/\log_E(P_0)=m_i/h\in\Z_p$.
Lemma~\ref{lem:quartic-uniform-reciprocal} and $e_i\le e_0$ imply
$C_i\equiv m_i/h\pmod {p^{e_i}}$. The remaining primary tests yield
$b_i\in\Z_p$ and $b_i\equiv n_i\pmod {p^{e_i}}$, as claimed.
\end{proof}

\begin{lemma}[Scalar height and the new anchor]
\label{lem:quartic-height19-anchor}
Every solution of the specified scalar batch and anchor has, with
$a=|s|$,
\[
 s\in2\cdot3^{2^{24}-1}\Z\setminus\{0\},\qquad
 h(\beta)\le H a^{19/32}<a/40,
\]
and
\begin{equation}\label{eq:quartic-cutoffs19}
 a>(24H)^{10},\qquad a>(40H)^{32/13},\qquad
 a>40,\qquad \log(40H)<a/40.
\end{equation}
Conversely this anchor can be imposed together with any fixed finite
list of formal divisibilities on $s$.
\end{lemma}
\begin{proof}
The exact anchor subgroup calculation in
Section~\ref{subsec:o8-anchor-subgroup}, with $25$ squarings, gives the
stated subgroup. Its two local computations are unchanged; the exponent
of the final power of three is now $2^{25}$.
In particular
\[
 \log_2 a>\tfrac32(2^{24}-1)=25165822.5.
\]
The binary height bound and $\log_2 24<5$, $\log_2 40<6$ give
\[
 10\log_2(24H)<20939710<25165822.5,
\]
\[
 \tfrac{32}{13}\log_2(40H)
 <\tfrac{32}{13}\,2093972<5154400<25165822.5.
\]
These prove the first two cutoffs and $a>40$.

Apply the same constant-left norm argument to the scalar batch. At each
odd contact prime of $\beta$ outside $S_*$ it excludes negative odd raw
valuation and gives the six assigned coefficient orders. The first four
make $J_0=[j_0]P_*$ and $S$ formal at every place over that prime. The
last two imply $v_p(s/j_0)\ge32$, hence $v_p(s)\ge32$. This reasoning
also applies at inert primes: both rational coefficients are tested and
the extension is unramified outside two. The retained scalar raw
polynomial and its small-prime correction detect every prime of
$M_5(\beta)$, including an odd infinite contact. Therefore
\[
 M_5(\beta)^{32}\mid s,
 \qquad h(\beta)\le H M_5(\beta)^{19}\le Ha^{19/32}.
\]
The domain hypothesis holds because $\beta=X^2+2$ is in none of the four
branch values: the only positive candidate $2401$ would require the
nonsquare integer $2399$ to be a rational square. The second cutoff
now gives $h(\beta)<a/40$.
Since $32/13>2$, that cutoff also gives $a>(40H)^2$.
For $H\ge1$, $40\log(40H)<(40H)^2$ follows at $H=1$ and then by
differentiating the difference. This proves the last bound.
Finally the displayed anchor subgroup contains arbitrarily divisible
nonzero integers, exactly as required in the base completeness argument.
\end{proof}

% BLOCK FIRST_CONTRADICTION.
\begin{lemma}[The Lucas comparison needs twenty-nine positions]
\label{lem:quartic-selective-lucas-growth}
The assignment in \eqref{eq:quartic-mode-partition}, together with
the proxy equation, implies
\[
 t^2-s(s-2)tz+(s-2)^2z^2=1,
 \qquad |t|>\exp(|s|-3).
\]
\end{lemma}
\begin{proof}
Put $a=|s|$, $B_L=\max(|s|,|t|,|z|)$ and
\[
 \Delta=t^2-s(s-2)tz+(s-2)^2z^2-1.
\]
Every finite position of the six coefficients of $x(S),x(T),x(Z)$,
as well as all five common positions, retains all three Lucas
comparisons. Lemma~\ref{lem:quartic-selective-local}
and the proxy equation therefore give $p^{40}\mid\Delta$ at every
odd raw-contact prime needed by these six coefficients. By the parity
argument in Lemma~\ref{lem:quartic-rational37}, each of their normalized
contact radicals has its fortieth power dividing $|\Delta|$ if
$\Delta\ne0$. Choose the point among $S,T,Z$ whose index has size $B_L$.
The index-height estimate \eqref{eq:o8-indexheight} and
\eqref{eq:quartic-direct-height} give
\[
 B_L^2-6<7H|\Delta|^{19/40}
       <7\cdot5^{19/40}HB_L^{19/10}<16HB_L^{19/10},
\]
where $|\Delta|<5B_L^4$. As $B_L,H\ge1$, this implies
$B_L^{1/10}<22H<24H$, hence $B_L<(24H)^{10}$.
It contradicts $B_L\ge a$ and \eqref{eq:quartic-cutoffs19}.
Thus $\Delta=0$. Since $z\ne0$ and $a>40$, Lemma~\ref{lem:o8-lucas}
gives the asserted bound on $|t|$ itself. The point $A$ and index $n$
have not been used in this first contradiction.
\end{proof}

% BLOCK BETA_ONLY_VARIANT: retain this block for the 29/8 endpoint.
\begin{lemma}[Soundness of the beta-only variant]
\label{lem:quartic-beta-only}
The beta-only assignment in \eqref{eq:quartic-mode-partition} implies
$\beta=n$ and $X\in\Z$.
\end{lemma}
\begin{proof}
Suppose $\beta\ne n$, write its difference in lowest terms, and put
$N=|\operatorname{num}(\beta-n)|\ge1$.
All thirty-seven positions retain a beta comparison. At every relevant
prime $p$, the local conclusion gives $v_p(\beta-n)\ge12$; reducedness
then shows $p\nmid\operatorname{den}(\beta-n)$ and $p^{12}\mid N$.
Every normalized radical for all eight coefficients is covered.
Set $B=\max(|n|,|s|,|t|,|z|)$. The upper and lower heights of the point
realizing $B$ give
\[
 B^2<7HN^{19/12}+6,
 \qquad N\le2\exp(h(\beta))B.
\]
The second estimate follows by writing $\beta=b/c$ in lowest terms:
$|b|,|c|\le\exp(h(\beta))$ and $|n|\le B$.
Since $7\cdot2^{19/12}+6<40$ and $H,B\ge1$,
\[
 B^{5/12}<40H\exp(19h(\beta)/12),
 \qquad
 \log B<\tfrac{12}{5}\log(40H)+\tfrac{19}{5}h(\beta)
         <\tfrac{31}{200}a<a-3.
\]
But $B\ge|t|$, contradicting
Lemma~\ref{lem:quartic-selective-lucas-growth}. Hence $\beta=n\in\Z$.
The equation $X^2=\beta-2$ forces a rational $X$ to have denominator
one, so $X\in\Z$. This second step uses the larger maximum $B$;
there is no unproved bound on $n$ in terms of $B_L$.
\end{proof}

% BLOCK SOURCE_THEOREM: the selected 29-full / 8-beta source.
\begin{theorem}[Arithmetic interpretation of the selective source]
\label{thm:quartic-source-arithmetic}
With the arithmetic inputs and source-version conventions of
Section~\ref{sec:degree8-original}, the selective system specified
above has rational witnesses at input $X\in\Q$ if and only if $X\in\Z$.
The assertion also holds for its compressed primary-chart presentation
described below. The original input coordinate is unchanged.
\end{theorem}
\begin{proof}
Soundness is Lemma~\ref{lem:quartic-beta-only}.
For completeness let $X\in\Z$, set $\beta=n=X^2+2$ and $A=[n]P_*$.
Choose the scalar sign and small-prime squareclass correction as in the
base construction. Choose $J_0$ sufficiently deep at its finite set of
scalar places and then choose a positive nonzero $s$ in
$2\cdot3^{2^{24}-1}\Z$ with all additionally required divisibilities.
Arbitrarily deep choices give the scalar order $32$ conditions, make the
normalized scalar leaves integral and the norm root a unit at the
assigned primes, and preserve the local squareclasses at $2,3,17$.
The pole theorem supplies the scalar witnesses.
Lemma~\ref{lem:o8-lucas} supplies positive $t,z$ with the required
index equation. Put $S=[s]P_*,T=[t]P_*,Z=[z]P_*$ and set
$(\xi,\tau,\zeta)=(s,t,z)$.

Now fix the eight rational abscissa coefficients and apply
Lemma~\ref{lem:quartic-rational37}. Let $\mathcal P$ be the set of
distinct selected ordinary primes, with the same orientation and
reduction orders $d_{p,1},d_{p,2}$ as in the base proof. Put
\[
 j=\operatorname{lcm}_{p\in\mathcal P}d_{p,2},\qquad J=[j]P_*.
\]
At a position assigned to $p$ choose a private integer $h_p$ satisfying
\[
 h_p\equiv0\pmod {d_{p,1}},\qquad
 h_p\equiv4\pmod {d_{p,2}},\qquad
 h_p\equiv4\pmod {p^{40+v_p(j)}}.
\]
The inherited same-prime conditions
$\gcd(d_{p,1},d_{p,2})=4$ and $p\nmid d_{p,1}d_{p,2}$ prove compatibility.
The factor $v_p(j)$ includes orders from every other selected prime;
no cross-prime coprimality is assumed. Avoid $h_p=0,4,-4$.
For the type at this position take only its required points
\[
 P_0=[h_p]P_*,\qquad P_i=[h_pn_i]P_*\quad(i\in I).
\]
All are nonidentity. If a primary and its target had equal abscissa,
they would be equal or opposite, forcing $h_p=4$ or $-4$ because
$G$ is torsion-free and the corresponding $n_i$ is nonzero. Thus every
retained affine subtraction guard holds. Their differences have indices
$(h_p-4)n_i$, with $n_0=1$.

The primary points are formal at the first place, and $J$ and all
required differences at the second. The displayed congruences and
Lemma~\ref{lem:quartic-uniform-reciprocal} give every retained test order;
in beta-only positions the common precision $40$ choice is stronger
than the required precision $12$. Increase the private congruence
precision when necessary to make the norm leaves as deep as required.
Repeated occurrences of a prime have the same orientation and otherwise
independent private witnesses. Choose the single rational $\theta$ last,
to meet all actual coefficientwise requirements across all positions
and comparator types. Finite Chinese remaindering supplies it.

Each norm tree retains its constant leaf on the all-left path. At the
assigned prime its root is a unit; the positive raw value times its
square has square reduced denominator and the three required local
squareclasses. The last part of Lemma~\ref{lem:quartic-rational37}
therefore supplies every remaining pole witness. This constructs a
solution with the given $X$. The compressed presentation is equivalent
by the following explicit maps.
\end{proof}

% BLOCK COMPRESSED: replaces the old five-point-only lemma and G15 circuit.
\begin{lemma}[Compressed charts and exact evaluation for each type]
\label{lem:quartic-compressed-primary}
For each required primary point use $u_i,v_i,f_i,x_i,p_i\in K$ with
\[
 f_i=u_i^2,\qquad v_i^2+df_i=1,\qquad
 p_i^2=x_i-Df_i,\qquad x_if_i=1.
\]
Put $G_i=G_{\lceil e_i/4\rceil-1}(Df_i^2)$, with the same convention
for $G_0$, and replace the primary ratios by
\[
 \Lambda p_0=G_0,\qquad B_i=C_ip_i,\qquad B_i\Lambda=G_i
 \quad(i\in I).
\]
This presentation has exactly the same rational solution projection as
the full charts and ratios for the specified comparator. Every external
point coordinate and the rational input $X$ are fixed.
\end{lemma}
\begin{proof}
From a full chart set $p_i=y_if_i$, $\Lambda=-L_{e_0}(P_0)=G_0/p_0$
and $B_i=C_ip_i$. The old equations give the four displayed chart
equations and $z_i=-1/p_i$, so the ratio equations have exactly this sign.
Conversely global nonvanishing of every $G_i$ makes $\Lambda,p_i$
nonzero. Also $x_if_i=1$ gives $f_i\ne0$, and
\[
 f_ip_i^2=1-Df_i^2=v_i^2(1+df_i)
\]
makes $v_i\ne0$. Restore
\[
 y_i=x_ip_i,\qquad z_i=-1/p_i,\qquad
 h_i=u_ip_i,\qquad w_i=u_ip_i/v_i,
 \qquad L_i=z_iG_i.
\]
These satisfy every equation of \eqref{eq:o8-chart}; the old common
logarithm inverse is $1/L_0$ and $C_i=L_i/L_0$.
In the subtraction equations replace each $y_i$ by $x_ip_i$, keeping
all their guards and all difference parameters. No other coordinate
has been removed. These forward and reverse maps are rational on the
specified nonzero loci and preserve each remaining equation.
\end{proof}

The finite evaluations used here can themselves be written quadratically.
For precision $40$ bind four field pairs
\[
 w=Df^2,\qquad w_2=w^2,\qquad w_3=w_2w,\qquad U_6,
\]
with the fourth pair defined by the quadratic equation
\[
 U_6=\frac{a_3+a_4w+a_5w_2+
                  w_3(a_6+a_7w+a_8w_2+a_9w_3)}{a_6}.
\]
The quadratic output expression is
\[
 G_9(w)=a_0+a_1w+a_2w_2+a_6w_3U_6,
 \qquad a_6=-\frac{418}{16575}\ne0.
\]
For precision $12$ bind only $w=Df^2$ and use
$G_2(w)=1+a_1w+a_2w^2$. Substitute these expressions directly into
$\Lambda p_0=G_0$ and $B_i\Lambda=G_i$. Every product is quadratic in
the bound rational coordinates. Expansion gives each required $a_j$
coefficient exactly. The scale is one; no finite-error estimate is used
as a substitute for these polynomial identities.

% BLOCK SOURCE_BOUNDARY: replace old source counts/pins only from final artifact.
The source defined by these equations is a new arithmetic specification.
In particular neither deletion of unused comparisons nor the changes to
the denominator logarithm and norm tree assert a coordinate-forgetting
map between arbitrary old and new witness tuples. The preceding theorem
proves both implications for the rational input projection. The explicit
maps in the compressed-chart lemma concern only those chart changes.
The base compiler remains construction machinery; the polynomial
presented below is the single selected quartic obtained from this source.
Its independent source check must verify the complete modules, the
ordered expression lists and position assignment, all nonzero guards,
the original affine height targets, scalar orders, anchor terminal
coordinate, and the full flattening. The subsequent quadratic lift and
certificate replay preserve the same input coordinate. Those finite
checks establish source and polynomial identities; the arithmetic
projection statement rests on the theorem above with the inherited
arithmetic inputs and their existing source-version qualifications.

% END BODY. BIBLIOGRAPHY INSERTS follow as comments; move them inside the
% existing thebibliography environment, retaining the existing cited keys.
% \bibitem{diamond-im} F. Diamond and J. Im,
% \emph{Modular forms and modular curves}, in Seminar on Fermat's Last Theorem,
% CMS Conference Proceedings 17 (1995), 39--133, \S12.5.
% \bibitem{baker-et-al-finiteness} M. H. Baker, E. Gonz\'alez-Jim\'enez,
% J. Gonz\'alez and B. Poonen, \emph{Finiteness results for modular curves
% of genus at least 2}, American Journal of Mathematics 127 (2005),
% 1325--1387, \S\S3.1, 3.2.1.
% \bibitem{mascot-moduli} N. Mascot, \emph{Moduli-friendly Eisenstein series
% over the p-adics and the computation of modular Galois representations},
% Research in Number Theory 8 (2022), article 48,
% \href{https://doi.org/10.1007/s40993-022-00329-6}{doi:10.1007/s40993-022-00329-6},
% \S2.1, (2.1.8).

